% Separate editorial construction for the calculus-specific algebraic completeness gap.
\section{The Free Boolean Algebra and Natural-Deduction Completeness}
\label{sec:p02-free-boolean-completeness}

The algebraic completeness argument needs more than a name for a
Lindenbaum algebra. Its operations must be well defined, its Boolean
identities proved, its universal property established, and membership
in the premise filter related to the actual deduction calculus.
This section supplies those constructions for the primitive-negation
classical natural deduction of \cref{sheet:nd-cpl}.
The earlier statements
\cref{thm:p02-arbitrary-prop-compactness-bpi,thm:p02-arbitrary-prop-completeness}
retain their separate cited standing and attribution
\parencite{HowardRubin1998}; the proofs here are distinct editorial results.

Work in classical \(\mathsf{ZF}\). The quotient, universal property and
filter correspondence use no choice principle. Only the specified
ultrafilter-extension step uses \(\mathsf{BPI}\).
The choice-free countable completeness proof of
\cref{thm:prop-completeness} is used only to obtain closed finite
propositional schemata. That proof constructs maximal consistent sets by
an enumeration; it does not use this quotient or arbitrary-language
completeness. Thus the dependency is not circular.

\subsection{The exact language and finite proof support}
\label{sec:p02-fbc-syntax}

\begin{hbconvention}{def:p02-fbc-language}{Arbitrary atoms, unchanged finitary calculus}
Let \(P\) be any set of atomic symbols, distinguished from connective and
punctuation tokens by disjoint tags if necessary.
\(\operatorname{Form}(P)\) has the grammar of
\cref{sec:prop-syntax-semantics-01}, with \(p\in P\) in place of \(p_n\):
the primitive constructors are exactly
\(\bottom,\neg,\land,\lor,\to\).
Retain the source abbreviations
\[
 \topv=\neg\bottom,\qquad
 \alpha\leftrightarrow\beta
       =(\alpha\to\beta)\land(\beta\to\alpha).
\]
All formulas are finite. In particular, \(P=\varnothing\) still permits
\(\bottom\) and its compound formulas.
The finite right-associated conjunction uses the declared sequence order;
its empty value is \(\topv\).

Write \(\Gamma\vdash_P\alpha\) for the relation of
\cref{def:p02-arbitrary-prop-language}.
For a derivation \(D\), let \(\operatorname{oa}(D)\) be the finite set
of formulas at its open assumption occurrences. The permitted set
\(\Gamma\) may be infinite; the condition is
\(\operatorname{oa}(D)\subseteq\Gamma\), not that every member of
\(\Gamma\) occurs.
Every discharge is nonvacuous. Copies have globally fresh injectively
renamed labels. The source's merging relabeling may then mark selected
open occurrences of one formula with one fresh label before a new
root discharge. No repetition rule is introduced.
\end{hbconvention}

The set of formulas can be constructed by finite stages: start with
the tagged atoms and \(\bottom\), close under the four nonnullary
constructors, and take the union of the stages indexed by
\(\mathbb N\). Formation is a set construction in ZF, not a
well-ordering of \(P\). Structural induction also shows that each
formula has only finitely many atomic occurrences.

\begin{hblemma}{thm:p02-fbc-finite-syntax}{proved=annex:p02.fbc.finite-syntax}{Finite support, discharge, and closed finite schemata}
Derivability is monotone in its permitted assumption set and has finite
support. A derivation whose premises are already in \(\Gamma\) can be
padded with an actual fresh assumption occurrence of \(\alpha\), while
retaining its conclusion. Consequently
\[
 \Gamma\cup\{\alpha\}\vdash_P\beta
       \quad\Longleftrightarrow\quad
 \Gamma\vdash_P\alpha\to\beta
\]
holds with the stipulated nonvacuous discharge convention.

A uniform formula substitution into a closed finite derivation preserves
every rule and every nonvacuous discharge. Hence every two-valued
tautological schema on a fixed finite list of source atoms has a
closed derivation after those atoms are replaced by arbitrary members
of \(\operatorname{Form}(P)\).
No enumeration of \(P\), BPI, or arbitrary-language completeness is
used to obtain these closed schemata.
\end{hblemma}

\subsection{Boolean-valued soundness before forming a quotient}
\label{sec:p02-fbc-boolean-soundness}

A Boolean algebra means the set-sized bounded distributive complemented
lattice of \cref{def:p01-bpi-boolean-algebra}; a trivial algebra is
allowed unless nontriviality is expressly required.
For an assignment \(f:P\to B\), define
\(\llbracket\alpha\rrbracket_f\) by structural recursion, interpreting
\(\bottom\) as \(0_B\), primitive negation as Boolean complement,
conjunction and disjunction as meet and join, and implication by
\[
 \llbracket\alpha\to\beta\rrbracket_f
       =\neg\llbracket\alpha\rrbracket_f
                         \vee\llbracket\beta\rrbracket_f .
\]
Thus \(\llbracket\topv\rrbracket_f=1_B\).
This is a specified algebraic interpretation, not a representation
of \(B\) by sets or two-valued points.

\begin{hblemma}{thm:p02-fbc-boolean-soundness}{proved=annex:p02.fbc.boolean-soundness}{Rule-by-rule Boolean-valued soundness}
For every finite \(\mathsf{ND}_{\mathsf{CPL}}\) derivation \(D\)
with conclusion \(\alpha\), every Boolean algebra \(B\), and every
assignment \(f:P\to B\),
\[
 \bigwedge_{\delta\in\operatorname{oa}(D)}
                   \llbracket\delta\rrbracket_f
                 \ \leq_B\ \llbracket\alpha\rrbracket_f .
\]
The empty meet is \(1_B\); repeated occurrences contribute an idempotent
factor, not an additional premise requirement.
In particular a closed proof has value \(1_B\), and a closed proof
of \(\alpha\leftrightarrow\beta\) implies
\(\llbracket\alpha\rrbracket_f=\llbracket\beta\rrbracket_f\).
The assertion includes RAA and trivial target algebras.
\end{hblemma}

\subsection{Provable equivalence and the Boolean quotient}
\label{sec:p02-fbc-quotient}

\begin{hbdefinition}{def:p02-fbc-quotient}{The syntactic Lindenbaum quotient}
For \(\alpha,\beta\in\operatorname{Form}(P)\), put
\[
 \alpha\sim_P\beta
   \quad\text{iff}\quad
 \varnothing\vdash_P\alpha\leftrightarrow\beta,
 \qquad
 [\alpha]_P=\{\beta\in\operatorname{Form}(P):\beta\sim_P\alpha\}.
\]
The intended algebra \(L_P\) has underlying set
\(\{[\alpha]_P:\alpha\in\operatorname{Form}(P)\}\) and operations
\[
\begin{gathered}
 [\alpha]_P\wedge[\beta]_P=[\alpha\land\beta]_P,\qquad
 [\alpha]_P\vee[\beta]_P=[\alpha\lor\beta]_P,\\
 \neg[\alpha]_P=[\neg\alpha]_P,\qquad
 0_L=[\bottom]_P,\quad 1_L=[\topv]_P .
\end{gathered}
\]
Here connectives inside a class are formula constructors; operations
on classes are Boolean operations. The next theorem establishes their
well-definedness before they are used.
\end{hbdefinition}

\begin{hbtheorem}{thm:p02-fbc-quotient}{proved=annex:p02.fbc.quotient}{The syntactic quotient is a nontrivial Boolean algebra}
The relation \(\sim_P\) is an equivalence relation and a congruence
for every primitive constructor. The displayed operations are
single-valued independently of representatives and make \(L_P\)
a nontrivial Boolean algebra.
It satisfies commutativity, associativity, idempotence, absorption,
both bounds, both distributive laws and both complement laws.
Double complementation and both De Morgan laws also hold.
Moreover
\[
 [\alpha\to\beta]_P=\neg[\alpha]_P\vee[\beta]_P .
\]
No quotient representative selector or arbitrary-language completeness
is assumed.
\end{hbtheorem}

\begin{hbtheorem}{thm:p02-fbc-free}{proved=annex:p02.fbc.free}{Free universal property and the syntactic order}
Define \(\eta_P:P\to L_P\) by \(\eta_P(p)=[p]_P\).
For every Boolean algebra \(B\), including a trivial one, and every
map \(f:P\to B\), there is exactly one Boolean homomorphism
\(h_f:L_P\to B\) such that \(h_f\eta_P=f\).
It is given without selecting representatives by
\[
 h_f([\alpha]_P)=\llbracket\alpha\rrbracket_f .
\]
The map \(\eta_P\) is injective. When \(P=\varnothing\), \(L_P\)
has exactly the two distinct elements \(0_L,1_L\) and still has
the stated free property.
For all formulas,
\[
 [\alpha]_P\leq_L[\beta]_P
       \quad\Longleftrightarrow\quad
 \varnothing\vdash_P\alpha\to\beta .
\]
\end{hbtheorem}

\subsection{The premise filter and exact derivability}
\label{sec:p02-fbc-premise-filter}

\begin{hbdefinition}{def:p02-fbc-premise-filter}{Filter generated by premises}
For any set \(\Gamma\subseteq\operatorname{Form}(P)\), put
\[
 F_\Gamma=
 \left\{b\in L_P:
 \begin{array}{l}
 \text{for some }n\in\mathbb N\text{ and a sequence }
       \gamma_1,\ldots,\gamma_n\in\Gamma,\\
 {}[\gamma_1]_P\wedge\cdots\wedge[\gamma_n]_P\leq_L b
 \end{array}\right\}.
\]
An empty sequence contributes \(1_L\). Finite repetitions do not
change its meet. This is a definition by finite witnesses, not a
choice of a witness for every possible member.
\end{hbdefinition}

\begin{hbtheorem}{thm:p02-fbc-filter}{proved=annex:p02.fbc.filter}{Generated filter, deduction, and inconsistency}
\(F_\Gamma\) is the least Boolean filter containing every
\([\gamma]_P\), \(\gamma\in\Gamma\), and
\[
 [\varphi]_P\in F_\Gamma
             \quad\Longleftrightarrow\quad
 \Gamma\vdash_P\varphi .
\]
Consequently \(F_\Gamma\) is proper iff
\(\Gamma\not\vdash_P\bottom\); if it is improper it is all of \(L_P\).
For \(\Gamma=\varnothing\), \(F_\Gamma=\{1_L\}\).
Both directions retain finite proof support, including the case
where the witnessing premise conjunction is empty.
\end{hbtheorem}

\subsection{Avoidance, the truth lemma, and completeness}
\label{sec:p02-fbc-completeness}

\begin{hbtheorem}{thm:p02-fbc-separation}{proved=annex:p02.fbc.separation}{One explicit BPI extension supplies a countervaluation}
Assume \(\mathsf{BPI}\) and \(\Gamma\not\vdash_P\varphi\).
Put \(a=[\varphi]_P\) and form
\[
 H=\{b\in L_P:
          \text{some }f\in F_\Gamma\text{ has }f\wedge\neg a\leq_L b\}.
\]
This is a proper filter containing \(F_\Gamma\) and \(\neg a\).
Boolean filter extension from \cref{thm:p01-bpi-equivalence} gives
an ultrafilter \(U\supseteq H\), hence \(a\notin U\).
Define the total two-valued assignment
\[
 v_U(p)=\chi_U([p]_P),\qquad
 \chi_U(b)=
 \begin{cases}1,&b\in U,\\0,&b\notin U.\end{cases}
\]
For every formula, its recursive truth value satisfies
\[
 \widehat v_U(\alpha)=\chi_U([\alpha]_P).
\]
Thus \(v_U\) satisfies every member of \(\Gamma\) and falsifies
\(\varphi\). The properness of \(H\) is proved, not included as an
unstated hypothesis of an avoidance helper.
\end{hbtheorem}

\begin{hbtheorem}{thm:p02-fbc-completeness}{proved=annex:p02.fbc.completeness}{BPI-relative completeness and its exact choice calibration}
Over \(\mathsf{ZF}+\mathsf{BPI}\), for every set \(P\), every set
\(\Gamma\subseteq\operatorname{Form}(P)\), and every formula \(\varphi\),
\[
 \Gamma\models_P\varphi
       \quad\Longleftrightarrow\quad
 \Gamma\vdash_P\varphi .
\]
Every consequence has a finite assumption support. Every finitely
satisfiable set of formulas is satisfiable.
Over \(\mathsf{ZF}\), the scheme asserting this completeness for
all sets \(P\) is equivalent to \(\mathsf{BPI}\), using the
already constructed unrestricted compactness bridge
\cref{thm:p01-bpi-propositional-compactness}.
The fixed countable-alphabet theorem remains choice-free and is not
strengthened merely by allowing arbitrary formula substitutions.
\end{hbtheorem}

\begin{hbexample}{def:p02-fbc-empty-boundaries}{The empty alphabet and an inconsistent filter}
With no atoms, the unique empty assignment evaluates every formula;
the quotient still distinguishes \([\bottom]\) from \([\neg\bottom]\).
No proper filter exists in a trivial target Boolean algebra, but the
free universal property into that algebra is valid.
If \(\Gamma\) contains both \(p\) and \(\neg p\), its two classes have
zero meet, so \(F_\Gamma=L_P\); there is no ultrafilter extending that
improper filter. If \(\Gamma\) is empty, its filter is \(\{1_L\}\);
an avoided conclusion must therefore have class different from \(1_L\).
\end{hbexample}

\begin{proofboundary}{def:p02-fbc-source-boundary}{What the construction uses}
The quotient and free property are syntactic ZF results, supported by
the actual nonvacuous-discharge calculus, closed finite schemata and
Boolean-valued soundness. They do not assume Boolean representability.
The one unrestricted extension in the countermodel construction uses BPI,
not Zorn's lemma, AC or a well-order of the atoms or formulas.
The two earlier cited statements and any separately allocated original
source presentations retain their distinct provenance and standing.
The decision procedures of \cref{thm:p02-decidability} are a different
kind of result: they act on finite encoded inputs over coded atoms.
They neither decide satisfiability of infinite premise sets over an
arbitrary set of atoms, which may admit no coding at all, nor construct
the valuation \(v_U\) supplied through \(\mathsf{BPI}\); the exact
statements and proofs are in \cref{sec:p02-w02-p02-017-finite-input-scope}.
\end{proofboundary}

\index{Boolean algebra!free syntactic construction}
\index{Lindenbaum algebra}\index{completeness!Boolean algebra method}
\index{Boolean Prime Ideal principle!propositional completeness}
